% ============================================================
%  Paper 07: Fine-Structure Constant from Spectral Volume of the D*-Sphere
%  Target: RevTeX 4-2 Compilation (SDGFT Series)
% ============================================================
\documentclass[amsmath,amssymb,aps,prd,twocolumn,superscriptaddress,nofootinbib]{revtex4-2}

\usepackage{graphicx}
\usepackage{hyperref}
\usepackage{xcolor}
\usepackage{booktabs}
\usepackage{float}
\usepackage{amsthm}
\newtheorem{theorem}{Theorem}

% ── SDGFT shared macros ──────────────────────────────────────
\usepackage{sdgft-macros}

% ── Document ─────────────────────────────────────────────────
\begin{document}

\preprint{SDGFT-2026-07}

\title{Fine-Structure Constant from Spectral Volume of the D*-Sphere}

\author{David A. Besemer}
\email{david.besemer@sdgft.org}
\affiliation{Independent Researcher}

\date{\today}

\begin{abstract}
We derive the inverse fine-structure constant alpha_em^(-1) as the volume of a D*-dimensional sphere, corrected by the elementary lattice tension. Using the exact tree-level and fixed-point values of D*, we find alpha_em^(-1) ~ 136.96, in excellent agreement with the high-energy value of the coupling constant at the Planck scale.
\end{abstract}

\maketitle

\section{Introduction}
We calculate the electromagnetic coupling constant at the Planck scale. The formulation uses the geometric volume of a sphere in fractional dimensions.

\section{Theoretical Formulation}
Here we formulate the core mathematical relations governing the system:
\begin{equation}
  \aem^{-1} = 2\pi (\Dstar)^3 + \dd\,\Dstar
\end{equation}
\begin{equation}
  \aem^{-1}(\Dstartree) = 2\pi \left(\frac{67}{24}\right)^3 + \frac{1}{24}\left(\frac{67}{24}\right) \approx 136.955
\end{equation}
\section{Geometric Volume Calculation}
The formula combines the spatial volume 2*pi*(D*)^3 with a boundary tension term delta*D*, representing the energy density of the lattice boundary.

\section{Conclusion}
Electromagnetism emerges as a boundary effect of the fractal spacetime manifold, with a coupling fixed by the spectral dimension.



\begin{thebibliography}{99}
\bibitem{Goldstein_Paper01} D. A. Besemer, ``Axiomatic Foundations of SDGFT,'' SDGFT-2026-01 (in preparation).
\bibitem{Goldstein_Paper04} D. A. Besemer, ``Effective Spectral Dimension and the Banach Fixed-Point Theorem in SDGFT,'' SDGFT-2026-04.
\bibitem{Goldstein_Code} D. A. Besemer, \texttt{sdgft} Python package, \url{https://github.com/david-besemer/sdgft} (2026).
\end{thebibliography}

\end{document}
